\clearpage
\section{From packet calculus to field dynamics}
The construction developed here has several exact outputs. It reconstructs operators from their transfer blocks. It converts skew-adjoint state dynamics into antisymmetric packet currents. It carries finite balance equations through partitions and moving frames. It identifies the additional information required by mixed states and open dynamics. It also gives precise versions of reconstruction beyond a finite atlas, where topology and domain membership become part of the statement.

The distinction between a supplied evolution and an existence theorem matters throughout the series. Differentiating a trajectory that already satisfies a state equation establishes the associated balance law. Constructing an evolution from a generator requires separate hypotheses and an existence argument. Identifying the generator of a previously constructed group is another specific task. Keeping these steps separate allows a later paper to reuse the balance identities without silently importing an unproved dynamics theorem.

The free Dirac monograph gives one concrete continuation. Fourier modes resolve the spinor field. Each mode carries a finite matrix, and a unitary Fourier map reconstructs the physical field. Finite packet consistency extends by Hilbert-space stability, while the maximal operator domain is recovered from weighted square summability. The Fourier construction provides the self-adjoint physical Hamiltonian. The following generator chapter proves its spectral-evolution identification at the level of mathematical analysis. Formal certification of that identification is recorded separately from the operator construction.

Gauge covariance and geometric refinement then extend the transport setting. Its task is to carry ordered transport through edge subdivision and cubical reconstruction while controlling the resulting defects. The moving-frame and transfer-block formulas in this volume supply part of its language. A regular prescribed connection, a coupled field evolution, and a constrained dynamical gauge theory each require their own precise mathematical setting.

The remaining parts of this book develop adic refinement, constrained gravity, and cosmological transport. Their use of the same resolution language does not eliminate their separate analytic and physical obligations. It makes those obligations comparable: specify the state space, identify the maps between resolutions, state the conserved or transported quantity, and exhibit the residual that prevents a stronger conclusion.

\begin{readerroadmap}[title={A reusable reading discipline}]
For each later result, locate the resolution, the transfer rule, the reconstruction statement, and the hypotheses that justify differentiation or passage to a limit. When a step fails, retain the explicit obstruction. It becomes part of the mathematical record and guides the next construction.
\end{readerroadmap}

The same reading discipline applies across the book: exact algebra, analytic reconstruction, and physical interpretation must each carry their own hypotheses.
